\chapter{Known Limitations}
\label{ch:known-limitations}

\textit{Every research system has limitations.  This chapter documents
seven known limitations of \sysname{}, with evidence, impact assessment,
and potential mitigations for each.  These are not bugs to fix---they are
boundaries of the current design that constrain what conclusions can be
drawn from the system and its evaluation.}

\begin{keyinsight}
Every limitation listed here is a potential research contribution for a
future student.  Addressing even one of these would strengthen the system
significantly and could constitute a publishable result.
\end{keyinsight}


% ─────────────────────────────────────────────────────────────────────
\section{Synthetic-Only Evaluation}
\label{sec:lim-synthetic}

\noindent\textbf{Description.}
All evaluation data comes from 200 synthetically generated conversations.
No real developer--AI conversations were used for evaluation.

\medskip\noindent\textbf{Evidence.}
The conversations in \codefile{evaluation/generate\_conversations.py} are
hand-authored templates that follow predictable patterns: a developer asks
about a technology choice, the assistant presents options, and they
converge on a decision.  Real conversations are messier---they include
tangents, incomplete thoughts, code blocks, and implicit decisions that
are never explicitly stated.

\medskip\noindent\textbf{Impact.}
Reported metrics (B-cubed F1 = 0.979) are likely \emph{optimistic}.
Entity mentions in synthetic data use well-known surface forms that the
canonical dictionary was designed to handle.  Real conversations may
contain novel abbreviations, domain-specific jargon, or typos that fall
outside the dictionary's coverage.

\medskip\noindent\textbf{Potential mitigation.}
Collect real conversations from open-source projects that use AI coding
assistants.  Several open-source teams publish their Claude Code or
Copilot Chat logs.  These could be anonymized and used as an evaluation
corpus (see Chapter~\ref{ch:future-directions},
Section~\ref{sec:future-semester}).


% ─────────────────────────────────────────────────────────────────────
\section{No User Study}
\label{sec:lim-no-user-study}

\noindent\textbf{Description.}
\sysname{}'s core value proposition---that a decision knowledge graph
improves developer decision-making---has not been validated with actual
users.

\medskip\noindent\textbf{Evidence.}
The system has been evaluated purely on technical metrics: entity
resolution accuracy, extraction success rates, and graph topology
statistics.  There is no data on whether developers who use \sysname{}
make better, faster, or more consistent decisions compared to developers
who do not.

\medskip\noindent\textbf{Impact.}
Without user validation, the system may solve a problem that does not
exist in practice, or solve it in a way that does not match developer
workflows.  Technical correctness does not imply practical utility.

\medskip\noindent\textbf{Potential mitigation.}
Conduct a controlled user study with two groups: one using \sysname{} and
one without.  Measure decision consistency, time to find relevant prior
decisions, and subjective usefulness ratings.


% ─────────────────────────────────────────────────────────────────────
\section{Weak Baselines}
\label{sec:lim-weak-baselines}

\noindent\textbf{Description.}
The entity resolution pipeline is not compared against strong published
baselines such as BLINK, EntQA, or fine-tuned BERT-based entity linkers.

\medskip\noindent\textbf{Evidence.}
The current evaluation compares the full pipeline against ablated versions
of itself (removing one stage at a time).  While this demonstrates the
value of each pipeline stage, it does not show whether the pipeline as a
whole is competitive with state-of-the-art entity linking systems.

\medskip\noindent\textbf{Impact.}
The high B-cubed F1 (0.979) may reflect the relative simplicity of the
synthetic evaluation data rather than genuine superiority of the approach.
A strong baseline would contextualize this number and reveal whether the
7-stage cascade is actually necessary.

\medskip\noindent\textbf{Potential mitigation.}
Implement at least one neural baseline (e.g.\ fine-tuned BERT for entity
matching) and one traditional baseline (e.g.\ TF-IDF cosine similarity,
as described in Exercise~\ref{sec:ex-new-baseline}).  Compare all systems
on the same held-out test set.


% ─────────────────────────────────────────────────────────────────────
\section{Two Dormant Stages}
\label{sec:lim-dormant-stages}

\noindent\textbf{Description.}
Two of the seven entity resolution stages---alias search (stage~4) and
embedding similarity (stage~6)---rarely trigger during evaluation.

\medskip\noindent\textbf{Evidence.}
The pipeline benchmark shows that the vast majority of mentions are
resolved by exact match (stage~2), canonical lookup (stage~3), or fuzzy
matching (stage~5).  Alias search and embedding similarity contribute
marginal improvements.  The evaluation data does not generate enough
hard cases to exercise these stages meaningfully.

\medskip\noindent\textbf{Impact.}
These stages add code complexity and runtime cost (embedding lookups
require vector similarity search) without demonstrable benefit on the
current evaluation data.  They may be valuable on harder, real-world data,
but this has not been shown.

\medskip\noindent\textbf{Potential mitigation.}
Design a targeted experiment with adversarial entity mentions that are
specifically crafted to bypass the first four stages.  If the dormant
stages resolve these correctly, their value is demonstrated.  If not,
consider simplifying the pipeline (see
Chapter~\ref{ch:future-directions}).


% ─────────────────────────────────────────────────────────────────────
\section{Single-Project Scope}
\label{sec:lim-single-project}

\noindent\textbf{Description.}
\sysname{} stores decisions within a single project context.  There is no
mechanism for cross-project knowledge transfer or sharing.

\medskip\noindent\textbf{Evidence.}
The data model uses \ic{user\_id} and \ic{project\_id} filtering on all
queries.  Entities and decisions from one project are invisible to
another, even when they refer to the same technologies.

\medskip\noindent\textbf{Impact.}
An organization using PostgreSQL across 10 projects would build 10
separate, redundant knowledge graphs.  Decisions made in one project
(e.g.\ ``we switched from REST to GraphQL because of N+1 query
problems'') cannot inform similar decisions in other projects.

\medskip\noindent\textbf{Potential mitigation.}
Implement cross-project entity alignment, where entities with the same
canonical name across projects are linked at the graph level.  This would
enable queries like ``what has our organization decided about
PostgreSQL?'' across all projects.


% ─────────────────────────────────────────────────────────────────────
\section{LLM Dependency}
\label{sec:lim-llm-dependency}

\noindent\textbf{Description.}
Decision extraction quality is tightly coupled to the specific LLM used
(NVIDIA Nemotron).  Results may differ substantially with a different
model.

\medskip\noindent\textbf{Evidence.}
The reproducibility analysis (Chapter~\ref{ch:evaluation-methodology},
Section~\ref{sec:eval-reproducibility}) shows that even the \emph{same}
model produces different results across runs: 364 decisions in Run~1 vs.\
383 in Run~2.  Switching to a different model (e.g.\ Claude, GPT-4) would
introduce additional variation from differences in instruction-following
behavior, JSON formatting reliability, and domain knowledge.

\medskip\noindent\textbf{Impact.}
The system's extraction quality is not a property of the pipeline
alone---it is a property of the pipeline \emph{plus} the specific model.
Users who substitute a weaker model may see significant quality
degradation.  Users who substitute a stronger model may see improvements
that the current evaluation does not predict.

\medskip\noindent\textbf{Potential mitigation.}
Run the evaluation with at least 3 different LLMs (e.g.\ Nemotron,
Claude, GPT-4) and report per-model metrics.  This would disentangle
pipeline quality from model quality and help users choose an appropriate
model for their use case.


% ─────────────────────────────────────────────────────────────────────
\section{Scale Untested}
\label{sec:lim-scale}

\noindent\textbf{Description.}
The system has been tested with approximately 200 conversations, 383
decisions, and 847 entities.  Performance at larger scales is unknown.

\medskip\noindent\textbf{Evidence.}
All evaluation data fits comfortably in memory.  The Neo4j graph has
$\sim$1,230 nodes and $\sim$1,271 relationships---well within the range
where brute-force operations are fast.  Entity resolution performs linear
scans in several stages.  The embedding similarity stage uses exact cosine
similarity, not approximate nearest neighbors.

\medskip\noindent\textbf{Impact.}
A production deployment with thousands of conversations and tens of
thousands of entities may encounter:

\begin{itemize}
    \item Slow entity resolution due to linear scans in the fuzzy matching
          stage.
    \item High memory usage for the embedding index.
    \item Slow graph traversal in the GraphRAG pipeline as subgraph
          expansion retrieves increasingly large neighborhoods.
    \item Neo4j query performance degradation without proper indexing.
\end{itemize}

\medskip\noindent\textbf{Potential mitigation.}
Profile the system at 10$\times$ and 100$\times$ current scale.  Replace
exact nearest-neighbor search with approximate nearest neighbors (ANN)
using FAISS or pgvector HNSW indexes.  Add query result pagination and
graph sampling for large neighborhoods (see
Chapter~\ref{ch:future-directions}, Section~\ref{sec:future-thesis}).
